\documentclass[onecolumn, draftclsnofoot,10pt, compsoc]{IEEEtran}
\usepackage{graphicx}
\usepackage{url}
\usepackage{setspace}
\usepackage{listings}

\usepackage{geometry}
\geometry{textheight=9.5in, textwidth=7in}

% 1. Fill in these details
\def \CapstoneTeamName{Manovacumeter Team}
\def \CapstoneTeamNumber{ 20}
\def \GroupMemberOne{Cade Raichart}
\def \GroupMemberTwo{Daniel Kato}
\def \GroupMemberThree{Nathan Shepherd}
\def \CapstoneProjectName{ISS Barometer App }
\def \CapstoneSponsorCompany{NASA}
\def \CapstoneSponsorPerson{Don Pettit}

% 2. Uncomment the appropriate line below so that the document type works
\def \DocType{Progress Report}

\newcommand{\NameSigPair}[1]{\par
\makebox[2.75in][r]{#1} \hfil 	\makebox[3.25in]{\makebox[2.25in]{\hrulefill} \hfill		\makebox[.75in]{\hrulefill}}
\par\vspace{-12pt} \textit{\tiny\noindent
\makebox[2.75in]{} \hfil		\makebox[3.25in]{\makebox[2.25in][r]{Signature} \hfill	\makebox[.75in][r]{Date}}}}
% 3. If the document is not to be signed, uncomment the RENEWcommand below
%\renewcommand{\NameSigPair}[1]{#1}

%%%%%%%%%%%%%%%%%%%%%%%%%%%%%%%%%%%%%%%
\begin{document}
\begin{titlepage}
    \pagenumbering{gobble}
    \begin{singlespace}
    	\includegraphics[height=4cm]{coe_v_spot1}
        \hfill
        % 4. If you have a logo, use this includegraphics command to put it on the coversheet.
        %\includegraphics[height=4cm]{NASA}
        \par\vspace{.2in}
        \centering
        \scshape{
            \huge CS Capstone \DocType \par
            {\large\today}\par
            \vspace{.5in}
            \textbf{\Huge\CapstoneProjectName}\par
            \vfill
            {\large Prepared for}\par
            \Huge \CapstoneSponsorCompany\par
            \vspace{5pt}
            {\Large\NameSigPair{\CapstoneSponsorPerson}\par}
            {\large Prepared by }\par
            Group\CapstoneTeamNumber\par
            % 5. comment out the line below this one if you do not wish to name your team
            %\CapstoneTeamName\par
            \vspace{5pt}
            {\Large
                \NameSigPair{\GroupMemberOne}\par
                \NameSigPair{\GroupMemberTwo}\par
                \NameSigPair{\GroupMemberThree}\par
            }
            \vspace{20pt}
        }
        \begin{abstract}
        % 6. Fill in your abstract
        The International Space Station, the conglomerate work of multiple national space associations, including NASA, requires precision devices to monitor the many aspects of space life.
        One very important aspect for the success of life in near vacuum is sustained air pressure inside the modules.
        This air pressure is monitored by mechanical barometers, called Manovacometers, which must be visually monitored, and are no longer available.
        The \CapstoneProjectName seeks to provide a suitable replacement for the current barometer device aboard the International Space Station.
        This document will serve to introduce the \CapstoneProjectName, and will detail the problem, the solution, and the success metrics of the project.
        The sections include the Definition of the Problem, Definition of Project, and the Metrics of Success.
        \end{abstract}
    \end{singlespace}
\end{titlepage}
\newpage
\pagenumbering{arabic}
\tableofcontents
% 7. uncomment this (if applicable). Consider adding a page break.
%\listoffigures
%\listoftables
\clearpage

% 8. now you write!
\section{Project Purpose}
The purpose of this project is to create an application, the ISS Barometer, that will be used on the ISS to measure and track air pressure.
The crew members aboard the ISS currently rely on out of date mechanical barometers to take critical readings during leaks.
These leaks, caused by orbital debris, must be found and patched or the crew would need to evacuate.
Our application seeks to improve on the mechanical barometers by providing clear and accurate pressure readings through the iPads on-board barometer sensor.
The app will display the current pressure and changes in pressure, by way of on screen digits and a plot.
Eventually, in version 2, our app will also have the capability of calculating remaining time on vessel before evacuation.
This added functionality will further assist the crew members, allowing them to make quick decisions in potential life and death situations.

\section{Project Goals}
As mentioned above the application we are making will provide ISS crew members with an accurate and efficient pressure reading.
Our required features follow that purpose and are listed below:
\begin{itemize}
\item[V1:] Data Screen with initial pressure, current pressure, start/stop button, time stamp, pressure change (measured in sec / mmHg)
\item[V1:] Graph Screen with plot pressure as a function of time
\item[V1:] Additional settings menu that includes plot scale, scroll functionality, orientation, decimal point, as well as other other units of measure for pressure
\item[V1:] All items will be optimized for efficient reading and clear interface
\item[V2:] Graph will have pinch to zoom functionality, scaling options (to be set in settings)
\item[V2:] Graph will be able to be exported as a csv, excel readable file
\item[V2:] Graph can be compared to t-Res table, to measure time remaining
\item[O1:] Research can be done to identify low pressure accuracy and resistance
\end{itemize}
These features, with the exception of the optional item O1, function as our goals.
Along with these specific features, the application that we create must be efficient and clear, as decided by the users.
The application will be used in challenging situations and must be clear and easy to understand, so that the data can be conveyed quickly and without confusion.

\section{Current State of the Project}
% If anyone can think of a better way to state this, feel free to change the heading
Currently we have an outline of what parts must be included in the app as well as the outline for when the first version of our app must be completed. We have also created our first prototype. This reads from the barometer and displays the current pressure in mm/Hg. It also calculates the dp/dt and displays in below the current pressure.

\section{Project Problems}
We chose to use swift and xcode because they allows for quick development and prototyping.
This said we ran into many problems with that. For one Swift only works on Linux and MacOS.
One of our group members did not have either at the time. They fixed this with installing a linux OS.
Another problem was installing the swift libraries and connecting them with the xcode application.
Linux needed many packages to be downloaded for the library to work correctly, this took many hours of troubleshooting as it was not straightforward.
Connecting swift to xcode on the MacOS was another problem as well.
Our last problem with swift and xcode is we have no experience using either.
We had and still have to continually work on improving our knowledge of these systems using online videos, guides and other resources.

For the first time we all experienced what it was like working with someone who is in the workplace.
We learned that our clients will not be getting back to us immediately and sometimes we have to wait for them reply.
This means we have to work and finish our work and send it to him way earlier then we had planned originally.
We also ran into the problems of finding the correct IEEE documentations.
Sometimes we found the document quick and easy while other times it hours and many students combined to find them.

\section{Interesting Code}
\begin{lstlisting}
  func startDisplayingPressureData() {
      self.altimeter.startRelativeAltitudeUpdates(to: OperationQueue.main, withHandler: {
        data, error in
          self.handlePressureReading(data: data!)
      })
  }

\end{lstlisting}

One of the main features of our application is its ability to read data from the iPad's built in barometer and display the data collected in different ways.
The way that data is collected from the barometer is quite interesting.
The process of reading data from the barometer requires a bit of processing power and time which can block other time sensitive tasks form being performed.
As a result, the measurement is done asynchronously through an object called the \texttt{OperationQueue}.
A method provided by the CoreMotion library called \texttt{startRelativeUpdates()} takes the \texttt{OperationQueue}'s \texttt{main} queue, and a callback function as arguement.
The function then appends a task to the queue to perform in the background, and when the task is complete, the result is passed into the callback function for further processing.\\

\section{Retrospective}
\begin{center}
	\begin{tabular}{| l | p{0.3\linewidth} | p{0.3\linewidth} | p{0.3\linewidth} |}
		\hline
		& Positives & Deltas & Actions \\ \hline
    Week 1 & Chose and submitted choices for our projects
      & No changes, just challenges in choosing
      & By justifying our choices, the preferences became clearer \\ \hline
		Week 2 & The group met, we had all chosen the ISS project as our first pick
      & We emailed Don Pettit, for a meeting, and started our Problem Statements
      & We were able to meet on Friday, Don had a compete picture of the app\\ \hline			% Put content between '&'s
		Week 3 & Finished problem statements and turned into both our client and teacher
 & Figure out how to share app with client, whom does not know the github system
 & step by step instructions on how to download the app \\ \hline
		Week 4 & Finalized Problem Statement, got client approval
      & Spelling of MANOVACUMETER was incorrect, among other minor edits
      & Changes were made, and the final Problem Statement was submitted \\ \hline
    Week 5 & Started Requirements Document rough draft & Need to start Requirements Document & Requirements Document will be started \\ \hline
		Week 6 & First prototype with two basic display of barometer numbers
  finished the Requirements Document & & \\ \hline
  	Week 7 & Completed Technology Review document rough draft & Need to complete Technology Review rough draft & A rough draft of the Technology Review will be created \\ \hline
		Week 8 & Started Technology Review & Need to Start Technology Review & Technology Review will be started \\ \hline
		Week 9 &Finished Technology document
Started working on design document
 & Could not find IEEE paper & Searched more and asked around\\ \hline
		Week 10 & Completed Design Document and Progress Report & Need to complete Design Document and Progress Report & Design Document and Progress Report will be completed \\ \hline

	\end{tabular}
\end{center}

\newpage
\bibliographystyle{IEEEtran}
\bibliography{bibliography}
\end{document}
